\chapter{Outlook and further Research Options}
\label{ch:otl}

\section{Quantum Effects}

Even though the calculation done in this thesis was a purely classical one,
    the found equilibrium positions and fields are also valid when modeling the particles 
    and fields as quantum degrees of freedom.
For the single particle case, one can exploit the fact that there are multiple equilibrium positions that whose locations can be tuned via the tweezers detuning. 
For some particular detuning values they even merge. Thus, there are opportunities to potentially observe tunneling between these.

Increasing the number of particles is expected to yield rich physics coming from the interparticle interaction induced by the electromagnetic field \cite{doi:10.1126/science.abp9941},
    which will be amplified due to the presence of the cavity.

\section{Generalizing the coordinates}
In this thesis, we only looked at systems that satisfy $k_c y_i = n\pi$.
It would be interesting to study the equilibrium positions when we relax this restriction
    and how this would affect the dynamics.

An interesting point to check would be if there exists a function $y_i(\vec z)$ that keeps the system in an equilibrium. 
This would remove a degree of freedom for each particle
    and return the system to $N$ conditions and $N$ degrees of freedom.
    
In the $N=2$ case, this would allow us to keep $l = const$ and vary $\Delta$ completely independently.

When doing those generalizations, 
    it would also be interesting to try to find a potential that only depends on the particle positions directly.
Currently,
    we took a look at $\grad V(\vec z, \alpha_c) \overset{!}{=} \vec 0$
    here we have $\alpha_c$ as an independent variable.
And this degree of freedom is only implicitly constrained by the particle positions $\vec z$.
It would provide meaningful insight if we could find a closed-form potential of $V = V(\vec z)$.
Since this would enable us to use the intuition that comes with a potential.


\section{Higher Dimensions}

The maximum number of particles we looked at was $2$.
For this amount, we may find that the solutions lie on a manifold,
    which has the same dimension as the number of particles.
It would be interesting to look at how this changes for higher $N$
    and if this pattern continues.

It would also be interesting to see if there are configurations where no solution exists
    and the system cannot be brought into an equilibrium.

Another thing to work on for higher dimensions is to verify if the found positions are even stable.
For the $1$ particle case, we found that all positions are stable,
    but this may not be a pattern that continues ad infinitum.
Finding an expression for $N$ particles 
    and checking the stability for higher $N$ would yield very interesting insight.


%\section{Continuous curves}
%
%Another interesting case to look at would be considering a wire in the field.
%If there exists a smooth curve of particles in configuration space, 
%    where the system is in an equilibrium.
%This corresponds to the formal limit of $N \rightarrow\infty$.
%Depending on the material used, this could represent either densely packed particles
%    or a wire in the field
%
%This would change the classical condition to a functional condition,
%    under this limit, we could look at this system using continuum Hamiltonian mechanics.
%

\section{Experimental verification}

Of course, this work is purely theoretical up until this point.
It would also be very interesting to see whether the predictions made here align with the experimental data 
    and if any effects were neglected here that would need to be taken into account.

